\section{Introduction}
\label{sec:introduction}

Reinforcement learning with verifiable rewards (RLVR) has become a standard
approach for improving language-model reasoning on tasks with automatically
checkable outcomes.  Yet a gain on the source task does not imply a gain on an
unseen target benchmark.  Depending on the source--target pair, post-training
may transfer positively, have no measurable effect, or degrade target
performance.  This uncertainty makes it difficult to choose RLVR data and
weakens claims that higher source-benchmark accuracy reflects broadly improved
reasoning.

Existing evaluations largely characterize transfer after training and often
describe task relatedness using broad domain labels or symmetric similarity.
These views do not directly answer the prospective question faced by a
practitioner: given a base policy, a source pool, and an unlabeled target task,
which source data is likely to transfer?  They also leave open whether a pass@1
gain expands the set of target problems the model can empirically solve or
merely concentrates probability on already accessible solutions.

\begin{figure*}[t]
    \centering
    \figureorplaceholder{figures/fig1_overview.pdf}{0.96\textwidth}{60mm}
    \caption{Problem setting and overview of \method.  A source dataset
    provides verifiable RLVR feedback, whereas the target contributes only
    unlabeled development prompts.  Reasoning Transition Graphs yield a
    directed source-to-target coverage score, which is used for prospective
    transfer prediction and fixed-budget source-data selection.}
    \label{fig:overview}
\end{figure*}

We formulate cross-benchmark RLVR generalization as a directed,
target-label-free transfer problem.  Our key hypothesis is that transfer
depends not only on semantic similarity, but on whether the source supplies
learnable reasoning structures demanded by the unresolved portion of the
target.  To test this hypothesis, we introduce \method, which represents
problems using Reasoning Transition Graphs and measures source-to-target
coverage with Directed Reasoning Coverage.  The score is explicitly
directional and conditioned on the base policy's source learnability and
target uncertainty.

We use this score in two ways.  First, we evaluate whether it predicts transfer
on held-out benchmarks and domains.  Second, \drcselect greedily chooses a
target-relevant and structurally diverse source subset under a fixed token
budget.  We additionally introduce a Bayesian capability-flow analysis that
separates empirical support expansion, sharpening, retention, and forgetting.

Our contributions are:
\begin{itemize}
    \item We formalize prospective, directed, target-label-free transfer for
    cross-benchmark RLVR and construct a matched-control transfer atlas.
    \item We propose Reasoning Transition Graphs and Directed Reasoning
    Coverage for predicting source-to-target transfer.
    \item We introduce a capability-flow analysis that distinguishes empirical
    support expansion from probability sharpening and forgetting.
    \item We propose \drcselect, a fixed-budget source-data selection method
    that uses no target answers or rewards.
\end{itemize}

